% =========================================================================
% PROTOCOLO EXPERIMENTAL E BENCHMARKS COMPUTACIONAIS
% =========================================================================
\chapter{Protocolo Experimental e Benchmarks Computacionais}\label{sec:benchmarks}

Para avaliar o desempenho prático dos algoritmos em relação às suas garantias assintóticas teóricas, estruturamos uma bateria de testes experimentais documentada em \code{Experimentos/}. Este capítulo detalha o protocolo metodológico, as instâncias de teste e os resultados comparativos obtidos.


\section{Rigor Metodológico e Ambiente de Teste}\label{sec:bench_metodologia}

O protocolo experimental fundamenta-se nas seguintes diretrizes:

\begin{enumerate}
    \item \textbf{Isolamento Estrito do Tempo de CPU:} O tempo de CPU é medido com \lstinline{std::chrono::high_resolution_clock} (C++23), cronometrando exclusivamente a execução de \lstinline{max_flow()} ou \lstinline{min_cost_flow()}, desconsiderando operações de entrada e saída ($I/O$) e impressão.
    
    \item \textbf{Parâmetros de Compilação e Otimização:} O código-fonte é compilado com GCC 13+ com suporte ao padrão ISO C++23, sob o nível de otimização \lstinline{-O2}. Adicionalmente, utiliza-se \lstinline{-Wl,-z,stack-size=268435456} para prevenir estouros de pilha em buscas profundas (DFS).
    
    \item \textbf{Tratamento Estatístico e Controle de Timeout:} Cada instância é executada de forma independente por $N = 5$ repetições em regime estacionário (\textit{steady state}), descartando o ciclo de aquecimento de cache (\textit{warm-up}). Registram-se a média e o desvio-padrão. Fixa-se o tempo limite de 60 segundos por instância; execuções que excedem esse limiar são interrompidas como \textit{Time Limit Exceeded} (TLE).
\end{enumerate}


\section{Famílias de Instâncias Padronizadas da Literatura (DIMACS)}\label{sec:bench_instancias}

Para submeter os algoritmos a diferentes desafios estruturais de rede, empregam-se as coleções do \textit{First DIMACS International Algorithm Implementation Challenge} (veja~\cite{Jensen2022, ahuja1993network}):

\begin{itemize}
    \item \textbf{Família DIMACS Max-Flow:}
    \begin{itemize}
        \item \textbf{Washington (RLG e Line):} Redes em malha retangular e redes lineares com gargalos induzidos para forçar caminhos aumentantes longos, desafiando a dependência do método de Ford-Fulkerson em relação às capacidades;
        \item \textbf{Genrmf e RMF:} Grafos estruturados em grade tridimensional formados por planos bidimensionais interconectados com capacidades aleatórias, onde a fonte reside no primeiro plano e o sumidouro no último. Essa família impõe gargalos de corte mínimo, avaliando o comportamento de métodos de pré-fluxo e caminhos aumentantes em redes com múltiplos cortes estreitos.
    \end{itemize}
    
    \item \textbf{Família DIMACS Min-Cost Flow:}
    \begin{itemize}
        \item \textbf{Netgen:} Gerador para instâncias de transporte e transbordo com capacidades e custos distribuídos uniformemente;
        \item \textbf{Grid (Toroidal e Planar):} Redes em grade bidimensional com conectividade espacial e custos correlacionados com distâncias euclidianas.
    \end{itemize}
    
    \item \textbf{Grafos Sintéticos e Aplicações Reais:}
    \begin{itemize}
        \item \textbf{Grafos Aleatórios Erdos-Rényi ($G(n, p)$):} Instâncias esparsas ($|A| \approx 4|V|$) e densas ($|A| \approx 0.1|V|^2$) para validação da sensibilidade à densidade de arestas;
        \item \textbf{Redes de Segmentação de Imagens:} Grafos derivados da modelagem descrita na Subseção~\ref{sec:segmentacao}, com vizinhanças 4-conectadas e capacidades ponderadas por gradiente de intensidade luminosa.
    \end{itemize}
\end{itemize}


\section{Métricas Operacionais e Certificação Cruzada de Corretude}\label{sec:bench_metricas}

Além do tempo de execução, a instrumentação computacional coleta métricas que refletem a dinâmica de cada algoritmo:

\begin{enumerate}
    \item \textbf{Contadores de Fases e Operações Elementares:}
    \begin{itemize}
        \item Em Ford-Fulkerson e Edmonds-Karp, quantificam-se a quantidade de caminhos aumentantes e o comprimento médio desses caminhos;
        \item No algoritmo de Dinic, contabilizam-se as fases de construção do digrafo de níveis e a quantidade de avanços e retrocessos de ponteiro na DFS de fluxo bloqueador;
        \item Na família Push-Relabel, monitoram-se separadamente as operações de empurrão (\textit{pushes}, saturantes e não saturantes) e de elevação (\textit{relabels}), evidenciando a eliminação de vértices desconectados do sumidouro pela heurística de lacuna (\textit{Gap Heuristic});
        \item No algoritmo \textit{Network Simplex}, registram-se o total de pivoteamentos de base e a proporção de pivôs degenerados (onde o fluxo do arco entrante é zero).
    \end{itemize}
    
    \item \textbf{Verificação Automática da Otimalidade e Consistência (Validação Cruzada):} O programa de testes verifica se todos os algoritmos produzem o mesmo valor ótimo:
    \begin{equation}
        |f^*_{\text{Edmonds-Karp}}| = |f^*_{\text{Dinic}}| = |f^*_{\text{Push-Relabel}}|, \quad \text{e} \quad z^*_{\text{CycleCanceling}} = z^*_{\text{SSP}} = z^*_{\text{NetworkSimplex}}
    \end{equation}
    Essa checagem cruzada previne discrepâncias de implementação ou instabilidades numéricas.
\end{enumerate}


\section{Quadro-Molde de Comparação dos Motores Algorítmicos}

A Tabela~\ref{tab:benchmarks} sintetiza o comportamento operacional esperado de cada algoritmo frente às diferentes famílias de redes.

\begin{table}[!ht]
\centering
\small
\setlength{\tabcolsep}{4.5pt}
\begin{tabular}{p{3.8cm} p{5.8cm} p{6.1cm}}
\toprule
\multicolumn{1}{c}{\textbf{Motor Algorítmico (C++)}} & \multicolumn{1}{c}{\textbf{Família de Melhor Desempenho}} & \multicolumn{1}{c}{\textbf{Cenário Crítico e Gargalo Empírico}} \\
\midrule
\multicolumn{3}{c}{\bannertext{\hyperref[sec:bench_instancias]{Motores de Fluxo Máximo (Coleções DIMACS Washington e RMF)}}} \\
\midrule
\hyperref[ap:ford_fulkerson]{Ford-Fulkerson} & Capacidades pequenas e caminhos curtos & Capacidades exponenciais e redes com ciclos \\
\hyperref[ap:edmonds_karp]{Edmonds-Karp} & Redes esparsas com poucos aumentos globais & Redes em camadas profundas (custo de BFSs sucessivas) \\
\hyperref[ap:dinic]{Dinic} & Redes unitárias, grafos bipartidos e redes densas & Gargalos múltiplos em níveis sucessivos da rede \\
\hyperref[ap:push_relabel]{Push-Relabel (FIFO)} & Redes densas com cortes mínimos próximos à fonte & Cadeias longas sem saturação imediata de excessos \\
\hyperref[ap:push_relabel_improved]{Push-Relabel (Gap)} & Grafos de visão (grid 2D/3D) e instâncias DIMACS RMF & Gargalos com desconexões sucessivas do sumidouro (eliminadas por Gap) \\
\midrule
\multicolumn{3}{c}{\bannertext{\hyperref[sec:bench_instancias]{Motores de Fluxo de Custo Mínimo (Coleções DIMACS Netgen e Grid)}}} \\
\midrule
\hyperref[ap:cycle_canceling]{Cycle Canceling} & Redes pequenas com poucos ciclos de custo negativo & Custos residuais elevados com múltiplos ciclos de $-1$ \\
\hyperref[ap:successive_shortest_dijkstra]{SSP (Dijkstra + $\pi$)} & Redes de transbordo com fluxo total moderado & Redes com fluxos massivos e capacidades unitárias \\
\hyperref[ap:network_simplex]{Network Simplex} & Grafos em grade, instâncias Netgen e redes logísticas & Degenerescência cíclica severa (controlada por Cunningham) \\
\bottomrule\end{tabular}
\caption{Quadro de análise comparativa dos motores algorítmicos e sensibilidade estrutural frente aos conjuntos de teste.}
\label{tab:benchmarks}
\end{table}

Os scripts no diretório \code{Experimentos/} automatizam a execução deste protocolo e geram as tabelas e gráficos apresentados a seguir.


\section{Resultados Experimentais: Fluxo Máximo}\label{sec:bench_resultados_maxflow}

A Tabela~\ref{tab:results_maxflow} apresenta os tempos médios de execução (5 repetições, $\bar{t}$ em ms) e o fluxo ótimo $f^*$ obtidos pelos algoritmos de fluxo máximo nas instâncias DIMACS e grafos sintéticos. O menor tempo por instância está destacado em \textbf{negrito}.

\IfFileExists{tabelas/table_max_flow.tex}{
\begin{table}[!ht]
\centering
\small
\setlength{\tabcolsep}{4pt}
\begin{tabular}{llrrrrl}
\toprule
\textbf{Inst\^ancia} & \textbf{Algoritmo} & $|V|$ & $|A|$ & $f^*$ & $\bar{t}$ (ms) & \textbf{Status} \\
\midrule
genrmf\_long.max & FordFulkerson & 1024 & 4080 & 40352 & 1179.08 & OK \\
genrmf\_long.max & EdmondsKarp & 1024 & 4080 & 40352 & 9.16 & OK \\
genrmf\_long.max & Dinic & 1024 & 4080 & 40352 & \textbf{2.91} & OK \\
genrmf\_long.max & PushRelabel & 1024 & 4080 & 40352 & 4.48 & OK \\
genrmf\_long.max & PushRelabelImproved & 1024 & 4080 & 40352 & 3.01 & OK \\
\midrule
genrmf\_medium.max & FordFulkerson & 1024 & 4544 & 130458 & 4872.33 & OK \\
genrmf\_medium.max & EdmondsKarp & 1024 & 4544 & 130458 & 6.89 & OK \\
genrmf\_medium.max & Dinic & 1024 & 4544 & 130458 & 2.48 & OK \\
genrmf\_medium.max & PushRelabel & 1024 & 4544 & 130458 & 114.81 & OK \\
genrmf\_medium.max & PushRelabelImproved & 1024 & 4544 & 130458 & \textbf{2.22} & OK \\
\midrule
genrmf\_small.max & FordFulkerson & 256 & 1008 & 38328 & 134.88 & OK \\
genrmf\_small.max & EdmondsKarp & 256 & 1008 & 38328 & 0.72 & OK \\
genrmf\_small.max & Dinic & 256 & 1008 & 38328 & \textbf{0.38} & OK \\
genrmf\_small.max & PushRelabel & 256 & 1008 & 38328 & 7.03 & OK \\
genrmf\_small.max & PushRelabelImproved & 256 & 1008 & 38328 & 0.50 & OK \\
\midrule
genrmf\_wide.max & FordFulkerson & 1024 & 4608 & 518574 & 16978.90 & OK \\
genrmf\_wide.max & EdmondsKarp & 1024 & 4608 & 518574 & 12.10 & OK \\
genrmf\_wide.max & Dinic & 1024 & 4608 & 518574 & 4.12 & OK \\
genrmf\_wide.max & PushRelabel & 1024 & 4608 & 518574 & 119.14 & OK \\
genrmf\_wide.max & PushRelabelImproved & 1024 & 4608 & 518574 & \textbf{2.88} & OK \\
\midrule
wash\_line\_100.max & FordFulkerson & 10002 & 39387 & -1 & --- & TLE \\
wash\_line\_100.max & EdmondsKarp & 10002 & 39387 & 154561846 & 8284.29 & OK \\
wash\_line\_100.max & Dinic & 10002 & 39387 & 154561846 & 67.04 & OK \\
wash\_line\_100.max & PushRelabel & 10002 & 39387 & 154561846 & 27832.40 & OK \\
wash\_line\_100.max & PushRelabelImproved & 10002 & 39387 & 154561846 & \textbf{64.41} & OK \\
\midrule
wash\_line\_50.max & FordFulkerson & 2502 & 9693 & -1 & --- & TLE \\
wash\_line\_50.max & EdmondsKarp & 2502 & 9693 & 77459801 & 244.92 & OK \\
wash\_line\_50.max & Dinic & 2502 & 9693 & 77459801 & \textbf{5.82} & OK \\
wash\_line\_50.max & PushRelabel & 2502 & 9693 & 77459801 & 824.14 & OK \\
wash\_line\_50.max & PushRelabelImproved & 2502 & 9693 & 77459801 & 6.14 & OK \\
\midrule
wash\_mesh\_16.max & FordFulkerson & 258 & 752 & 14283 & 27.03 & OK \\
wash\_mesh\_16.max & EdmondsKarp & 258 & 752 & 14283 & 2.39 & OK \\
wash\_mesh\_16.max & Dinic & 258 & 752 & 14283 & \textbf{0.38} & OK \\
wash\_mesh\_16.max & PushRelabel & 258 & 752 & 14283 & 1.66 & OK \\
wash\_mesh\_16.max & PushRelabelImproved & 258 & 752 & 14283 & 0.56 & OK \\
\midrule
wash\_mesh\_32.max & FordFulkerson & 1026 & 3040 & 28917 & 662.94 & OK \\
wash\_mesh\_32.max & EdmondsKarp & 1026 & 3040 & 28917 & 36.77 & OK \\
wash\_mesh\_32.max & Dinic & 1026 & 3040 & 28917 & \textbf{3.31} & OK \\
wash\_mesh\_32.max & PushRelabel & 1026 & 3040 & 28917 & 58.66 & OK \\
wash\_mesh\_32.max & PushRelabelImproved & 1026 & 3040 & 28917 & 4.30 & OK \\
\midrule
wash\_rlg\_16.max & FordFulkerson & 258 & 752 & 12100 & 73.07 & OK \\
wash\_rlg\_16.max & EdmondsKarp & 258 & 752 & 12100 & 2.47 & OK \\
wash\_rlg\_16.max & Dinic & 258 & 752 & 12100 & \textbf{0.47} & OK \\
wash\_rlg\_16.max & PushRelabel & 258 & 752 & 12100 & 3.60 & OK \\
wash\_rlg\_16.max & PushRelabelImproved & 258 & 752 & 12100 & 0.80 & OK \\
\midrule
wash\_rlg\_32.max & FordFulkerson & 1026 & 3040 & 21933 & 763.99 & OK \\
wash\_rlg\_32.max & EdmondsKarp & 1026 & 3040 & 21933 & 38.97 & OK \\
wash\_rlg\_32.max & Dinic & 1026 & 3040 & 21933 & \textbf{2.04} & OK \\
wash\_rlg\_32.max & PushRelabel & 1026 & 3040 & 21933 & 56.54 & OK \\
wash\_rlg\_32.max & PushRelabelImproved & 1026 & 3040 & 21933 & 3.52 & OK \\
\midrule
wash\_rlg\_64.max & FordFulkerson & 4098 & 12224 & 45921 & 7342.96 & OK \\
wash\_rlg\_64.max & EdmondsKarp & 4098 & 12224 & 45921 & 695.67 & OK \\
wash\_rlg\_64.max & Dinic & 4098 & 12224 & 45921 & \textbf{16.37} & OK \\
wash\_rlg\_64.max & PushRelabel & 4098 & 12224 & 45921 & 395.14 & OK \\
wash\_rlg\_64.max & PushRelabelImproved & 4098 & 12224 & 45921 & 22.11 & OK \\
\midrule
grid\_100x100.max & FordFulkerson & 10000 & 39600 & 54 & 23.24 & OK \\
grid\_100x100.max & EdmondsKarp & 10000 & 39600 & 54 & 24.47 & OK \\
grid\_100x100.max & Dinic & 10000 & 39600 & 54 & \textbf{2.90} & OK \\
grid\_100x100.max & PushRelabel & 10000 & 39600 & 54 & 9589.17 & OK \\
grid\_100x100.max & PushRelabelImproved & 10000 & 39600 & 54 & 35.14 & OK \\
\midrule
grid\_10x10.max & FordFulkerson & 100 & 360 & 129 & 0.47 & OK \\
grid\_10x10.max & EdmondsKarp & 100 & 360 & 129 & 0.12 & OK \\
grid\_10x10.max & Dinic & 100 & 360 & 129 & \textbf{0.08} & OK \\
grid\_10x10.max & PushRelabel & 100 & 360 & 129 & 0.18 & OK \\
grid\_10x10.max & PushRelabelImproved & 100 & 360 & 129 & 0.12 & OK \\
\midrule
grid\_50x50.max & FordFulkerson & 2500 & 9800 & 42 & 4.01 & OK \\
grid\_50x50.max & EdmondsKarp & 2500 & 9800 & 42 & 2.86 & OK \\
grid\_50x50.max & Dinic & 2500 & 9800 & 42 & \textbf{0.56} & OK \\
grid\_50x50.max & PushRelabel & 2500 & 9800 & 42 & 538.71 & OK \\
grid\_50x50.max & PushRelabelImproved & 2500 & 9800 & 42 & 9.23 & OK \\
\midrule
random\_dense\_100.max & FordFulkerson & 100 & 1000 & 4987 & 10.29 & OK \\
random\_dense\_100.max & EdmondsKarp & 100 & 1000 & 4987 & 0.37 & OK \\
random\_dense\_100.max & Dinic & 100 & 1000 & 4987 & \textbf{0.30} & OK \\
random\_dense\_100.max & PushRelabel & 100 & 1000 & 4987 & 3.68 & OK \\
random\_dense\_100.max & PushRelabelImproved & 100 & 1000 & 4987 & 0.38 & OK \\
\midrule
random\_dense\_500.max & FordFulkerson & 500 & 25000 & 25347 & 1233.91 & OK \\
random\_dense\_500.max & EdmondsKarp & 500 & 25000 & 25347 & 12.46 & OK \\
random\_dense\_500.max & Dinic & 500 & 25000 & 25347 & \textbf{2.99} & OK \\
random\_dense\_500.max & PushRelabel & 500 & 25000 & 25347 & 401.78 & OK \\
random\_dense\_500.max & PushRelabelImproved & 500 & 25000 & 25347 & 3.34 & OK \\
\midrule
random\_sparse\_100.max & FordFulkerson & 100 & 400 & 985 & 1.26 & OK \\
random\_sparse\_100.max & EdmondsKarp & 100 & 400 & 985 & 0.17 & OK \\
random\_sparse\_100.max & Dinic & 100 & 400 & 985 & 0.18 & OK \\
random\_sparse\_100.max & PushRelabel & 100 & 400 & 985 & 0.15 & OK \\
random\_sparse\_100.max & PushRelabelImproved & 100 & 400 & 985 & \textbf{0.13} & OK \\
\midrule
random\_sparse\_1000.max & FordFulkerson & 1000 & 4000 & 914 & 45.46 & OK \\
random\_sparse\_1000.max & EdmondsKarp & 1000 & 4000 & 914 & \textbf{0.59} & OK \\
random\_sparse\_1000.max & Dinic & 1000 & 4000 & 914 & 0.93 & OK \\
random\_sparse\_1000.max & PushRelabel & 1000 & 4000 & 914 & 127.62 & OK \\
random\_sparse\_1000.max & PushRelabelImproved & 1000 & 4000 & 914 & 1.99 & OK \\
\midrule
random\_sparse\_500.max & FordFulkerson & 500 & 2000 & 1572 & 30.84 & OK \\
random\_sparse\_500.max & EdmondsKarp & 500 & 2000 & 1572 & 0.85 & OK \\
random\_sparse\_500.max & Dinic & 500 & 2000 & 1572 & \textbf{0.53} & OK \\
random\_sparse\_500.max & PushRelabel & 500 & 2000 & 1572 & 31.48 & OK \\
random\_sparse\_500.max & PushRelabelImproved & 500 & 2000 & 1572 & 0.94 & OK \\
\midrule
worst\_chain\_100.max & FordFulkerson & 10002 & 10100 & 100 & \textbf{0.34} & OK \\
worst\_chain\_100.max & EdmondsKarp & 10002 & 10100 & 100 & 9.36 & OK \\
worst\_chain\_100.max & Dinic & 10002 & 10100 & 100 & 0.57 & OK \\
worst\_chain\_100.max & PushRelabel & 10002 & 10100 & 100 & 0.79 & OK \\
worst\_chain\_100.max & PushRelabelImproved & 10002 & 10100 & 100 & 0.86 & OK \\
\midrule
worst\_ff\_10.max & FordFulkerson & 22 & 41 & 2048 & 0.15 & OK \\
worst\_ff\_10.max & EdmondsKarp & 22 & 41 & 2048 & 0.07 & OK \\
worst\_ff\_10.max & Dinic & 22 & 41 & 2048 & 0.06 & OK \\
worst\_ff\_10.max & PushRelabel & 22 & 41 & 2048 & 0.06 & OK \\
worst\_ff\_10.max & PushRelabelImproved & 22 & 41 & 2048 & \textbf{0.05} & OK \\
\midrule
worst\_ff\_15.max & FordFulkerson & 32 & 61 & 65536 & 0.13 & OK \\
worst\_ff\_15.max & EdmondsKarp & 32 & 61 & 65536 & \textbf{0.04} & OK \\
worst\_ff\_15.max & Dinic & 32 & 61 & 65536 & 0.04 & OK \\
worst\_ff\_15.max & PushRelabel & 32 & 61 & 65536 & 0.04 & OK \\
worst\_ff\_15.max & PushRelabelImproved & 32 & 61 & 65536 & 0.04 & OK \\
\midrule
test\_small.max & FordFulkerson & 6 & 8 & 15 & 0.11 & OK \\
test\_small.max & EdmondsKarp & 6 & 8 & 15 & 0.06 & OK \\
test\_small.max & Dinic & 6 & 8 & 15 & \textbf{0.03} & OK \\
test\_small.max & PushRelabel & 6 & 8 & 15 & 0.04 & OK \\
test\_small.max & PushRelabelImproved & 6 & 8 & 15 & 0.06 & OK \\
\bottomrule
\end{tabular}
\caption{Resultados experimentais dos motores de fluxo m\'aximo.}
\label{tab:results_maxflow}
\end{table}

}{}

\IfFileExists{figuras/scalability_maxflow.pdf}{
\begin{figure}[!ht]
    \centering
    \includegraphics[width=0.92\textwidth]{figuras/scalability_maxflow}
    \caption{Escalabilidade dos algoritmos de fluxo máximo em função de $|V|$ (escala log-log).}
    \label{fig:scalability_maxflow}
\end{figure}
}{}

\IfFileExists{figuras/bars_maxflow.pdf}{
\begin{figure}[!ht]
    \centering
    \includegraphics[width=0.98\textwidth]{figuras/bars_maxflow}
    \caption{Comparação direta dos tempos médios de execução por instância (algoritmos de fluxo máximo).}
    \label{fig:bars_maxflow}
\end{figure}
}{}


\section{Resultados Experimentais: Fluxo de Custo Mínimo}\label{sec:bench_resultados_mincost}

A Tabela~\ref{tab:results_mincost} apresenta os resultados para os algoritmos de fluxo de custo mínimo nas instâncias Netgen da suíte DIMACS, reportando o custo ótimo $z^*$.

\IfFileExists{tabelas/table_min_cost.tex}{
\begin{table}[!ht]
\centering
\small
\setlength{\tabcolsep}{4pt}
\begin{tabular}{llrrrrrrl}
\toprule
\textbf{Inst\^ancia} & \textbf{Algoritmo} & $|V|$ & $|A|$ & $f^*$ & $z^*$ & $\bar{t}$ (ms) & \textbf{Status} \\
\midrule
test\_small.min & CycleCanceling & 4 & 5 & 6 & 24 & 0.00 & OK \\
test\_small.min & SuccessiveShortest & 4 & 5 & 6 & 24 & 0.00 & OK \\
test\_small.min & SuccessiveShortestDijkstra & 4 & 5 & 6 & 24 & \textbf{0.00} & OK \\
test\_small.min & NetworkSimplex & 4 & 5 & 6 & 24 & 0.00 & OK \\
\midrule
netgen\_01.min & CycleCanceling & 202 & 1508 & 100000 & 196587626 & 542.60 & OK \\
netgen\_01.min & SuccessiveShortest & 202 & 1508 & 100000 & 196587626 & \textbf{10.21} & OK \\
netgen\_01.min & SuccessiveShortestDijkstra & 202 & 1508 & 100000 & 196587626 & 13.08 & OK \\
netgen\_01.min & NetworkSimplex & 202 & 1508 & 100000 & 196587626 & 12.86 & OK \\
\midrule
netgen\_02.min & CycleCanceling & 202 & 1711 & 100000 & 194072029 & 605.62 & OK \\
netgen\_02.min & SuccessiveShortest & 202 & 1711 & 100000 & 194072029 & \textbf{9.91} & OK \\
netgen\_02.min & SuccessiveShortestDijkstra & 202 & 1711 & 100000 & 194072029 & 12.52 & OK \\
netgen\_02.min & NetworkSimplex & 202 & 1711 & 100000 & 194072029 & 12.40 & OK \\
\midrule
netgen\_03.min & CycleCanceling & 202 & 2200 & 100000 & 159442947 & 1061.77 & OK \\
netgen\_03.min & SuccessiveShortest & 202 & 2200 & 100000 & 159442947 & \textbf{14.82} & OK \\
netgen\_03.min & SuccessiveShortestDijkstra & 202 & 2200 & 100000 & 159442947 & 15.72 & OK \\
netgen\_03.min & NetworkSimplex & 202 & 2200 & 100000 & 159442947 & 14.89 & OK \\
\midrule
netgen\_04.min & CycleCanceling & 202 & 2400 & 100000 & 138936551 & 1189.21 & OK \\
netgen\_04.min & SuccessiveShortest & 202 & 2400 & 100000 & 138936551 & \textbf{15.17} & OK \\
netgen\_04.min & SuccessiveShortestDijkstra & 202 & 2400 & 100000 & 138936551 & 15.63 & OK \\
netgen\_04.min & NetworkSimplex & 202 & 2400 & 100000 & 138936551 & 15.42 & OK \\
\midrule
netgen\_05.min & CycleCanceling & 202 & 3100 & 100000 & 102950805 & 2115.92 & OK \\
netgen\_05.min & SuccessiveShortest & 202 & 3100 & 100000 & 102950805 & 24.69 & OK \\
netgen\_05.min & SuccessiveShortestDijkstra & 202 & 3100 & 100000 & 102950805 & 21.35 & OK \\
netgen\_05.min & NetworkSimplex & 202 & 3100 & 100000 & 102950805 & \textbf{20.98} & OK \\
\midrule
netgen\_06.min & CycleCanceling & 302 & 3474 & 150000 & 191968577 & 4143.53 & OK \\
netgen\_06.min & SuccessiveShortest & 302 & 3474 & 150000 & 191968577 & 43.23 & OK \\
netgen\_06.min & SuccessiveShortestDijkstra & 302 & 3474 & 150000 & 191968577 & \textbf{41.47} & OK \\
netgen\_06.min & NetworkSimplex & 302 & 3474 & 150000 & 191968577 & 43.27 & OK \\
\midrule
netgen\_07.min & CycleCanceling & 302 & 4819 & 150000 & 172742047 & 6674.34 & OK \\
netgen\_07.min & SuccessiveShortest & 302 & 4819 & 150000 & 172742047 & 59.52 & OK \\
netgen\_07.min & SuccessiveShortestDijkstra & 302 & 4819 & 150000 & 172742047 & \textbf{47.88} & OK \\
netgen\_07.min & NetworkSimplex & 302 & 4819 & 150000 & 172742047 & 51.20 & OK \\
\midrule
netgen\_08.min & CycleCanceling & 302 & 5469 & 150000 & 177575586 & 7856.27 & OK \\
netgen\_08.min & SuccessiveShortest & 302 & 5469 & 150000 & 177575586 & 66.33 & OK \\
netgen\_08.min & SuccessiveShortestDijkstra & 302 & 5469 & 150000 & 177575586 & \textbf{49.04} & OK \\
netgen\_08.min & NetworkSimplex & 302 & 5469 & 150000 & 177575586 & 54.70 & OK \\
\bottomrule
\end{tabular}
\caption{Resultados experimentais dos motores de fluxo de custo m\'inimo.}
\label{tab:results_mincost}
\end{table}

}{}

\IfFileExists{figuras/scalability_mincost.pdf}{
\begin{figure}[!ht]
    \centering
    \includegraphics[width=0.92\textwidth]{figuras/scalability_mincost}
    \caption{Escalabilidade dos algoritmos de fluxo de custo mínimo em função de $|V|$ (escala log-log).}
    \label{fig:scalability_mincost}
\end{figure}
}{}

\IfFileExists{figuras/bars_mincost.pdf}{
\begin{figure}[!ht]
    \centering
    \includegraphics[width=0.98\textwidth]{figuras/bars_mincost}
    \caption{Comparação direta dos tempos médios de execução por instância (algoritmos de fluxo de custo mínimo).}
    \label{fig:bars_mincost}
\end{figure}
}{}
